\documentclass[conference]{IEEEtran}
\IEEEoverridecommandlockouts

\usepackage[T1]{fontenc}
\usepackage[utf8]{inputenc}
\usepackage[brazil,shorthands=off]{babel}
\usepackage{cite}
\usepackage{amsmath,amssymb}
\usepackage{graphicx}
\usepackage{booktabs}
\usepackage{array}
\usepackage{xcolor}
\usepackage{listings}
\usepackage{url}
\usepackage{tikz}
\usetikzlibrary{automata,positioning,arrows.meta}

\definecolor{cinza}{gray}{0.45}
\definecolor{azul}{RGB}{20,70,140}

\lstdefinelanguage{Mare}{
  morekeywords={principal,funcao,seja,fixo,se,senao,enquanto,para,de,ate,passo,
    repita,retorne,mostre,leia,inteiro,real,logico,texto,verdadeiro,falso,e,ou,nao},
  sensitive=true,
  morecomment=[l]{\#},
  morecomment=[s]{\#*}{*\#},
  morestring=[b]"
}
\lstset{
  basicstyle=\ttfamily\footnotesize,
  keywordstyle=\color{azul}\bfseries,
  commentstyle=\color{cinza}\itshape,
  stringstyle=\color{teal},
  columns=fullflexible,
  keepspaces=true,
  showstringspaces=false,
  breaklines=true,
  frame=single,
  rulecolor=\color{black!30},
  numbers=left,
  numberstyle=\tiny\color{cinza},
  xleftmargin=1.6em,
  framexleftmargin=1.4em,
  aboveskip=4pt,
  belowskip=4pt,
  literate={á}{{\'a}}1 {ã}{{\~a}}1 {ç}{{\c{c}}}1 {é}{{\'e}}1 {ê}{{\^e}}1
           {í}{{\'i}}1 {ó}{{\'o}}1 {õ}{{\~o}}1 {ú}{{\'u}}1
}
\lstdefinestyle{ebnf}{
  language={}, numbers=none, xleftmargin=0.3em, framexleftmargin=0.2em,
  basicstyle=\ttfamily\scriptsize
}

\renewcommand{\lstlistingname}{Listagem}

\begin{document}

\title{Maré: Especificação e Implementação dos Analisadores Léxico e Sintático de uma Linguagem Imperativa Didática em Português}

\author{
\IEEEauthorblockN{Gustavo Teixeira Bittencourt de Oliveira, Christophe Abelem Pinho,\\
Edgar Klewert Da Costa Araujo, Elissandra Bernadett Abdon Nascimento, Adler Augustus de Castro Mota}
\IEEEauthorblockA{Curso de Engenharia da Computação -- Turma EC8MA\\
Centro Universitário do Estado do Pará (CESUPA)\\
Belém, PA, Brasil}
}

\maketitle

\begin{abstract}
Este artigo descreve a primeira etapa da construção de um compilador para a
Maré, uma linguagem de programação imperativa, estaticamente tipada e com
palavras-chave em português, projetada pela equipe. São apresentados a
especificação léxica da linguagem, com suas classes de lexemas, tokens e
expressões regulares, e uma gramática livre de contexto em notação EBNF que
cobre declarações de variáveis e constantes, o bloco principal, condicionais
encadeados, três formas de laço, funções com parâmetros tipados e expressões
aritméticas, relacionais e lógicas. O analisador léxico foi implementado em
Python, sem geradores automáticos, a partir de uma expressão regular mestre
com grupos nomeados. O analisador sintático é descendente recursivo, constrói
uma árvore sintática abstrata e recupera-se de erros em modo pânico. Um
programa de demonstração com 60 linhas é convertido em uma cadeia de 288 tokens
e aceito pelo analisador sintático, e 28 testes automatizados validam
precedência, associatividade e o diagnóstico de erros.
\end{abstract}

\begin{IEEEkeywords}
compiladores, análise léxica, análise sintática, expressões regulares,
gramática livre de contexto, parser descendente recursivo.
\end{IEEEkeywords}

% ======================================================================
\section{Introdução}

Um compilador traduz um programa escrito em uma linguagem-fonte para uma
representação equivalente em outra linguagem, e o faz em fases bem
delimitadas~\cite{aho2006}. As duas primeiras formam a \emph{vanguarda}
(\textit{front end}) do processo: a análise léxica agrupa os caracteres do
código em unidades com significado, os \emph{tokens}, e a análise sintática
verifica se a sequência desses tokens respeita a estrutura da linguagem,
produzindo uma árvore que servirá às fases seguintes~\cite{cooper2011}.

O objetivo deste trabalho é construir essas duas fases para uma linguagem
nova, a \textbf{Maré}, que servirá de protótipo ao longo do semestre. A
escolha por uma linguagem própria, em vez de um subconjunto de uma linguagem
existente, permitiu controlar cada decisão de projeto e justificá-la do ponto
de vista da teoria de linguagens formais. O nome faz referência ao fenômeno
das marés, característico da região amazônica, e também sugere o fluxo de
tokens que ``sobe'' do código-fonte até a árvore sintática.

As contribuições deste artigo são: (i) a especificação léxica e sintática da
Maré; (ii) um analisador léxico baseado em expressões regulares com detecção
de erros; (iii) um analisador sintático descendente recursivo com construção
de árvore sintática abstrata (AST) e recuperação de erros; e (iv) a validação
da implementação por meio de programas de exemplo e testes automatizados.

O restante do texto está organizado da seguinte forma: a Seção~\ref{sec:fund}
revisa os conceitos utilizados; a Seção~\ref{sec:lingua} apresenta a
linguagem; as Seções~\ref{sec:lex} e~\ref{sec:sint} detalham os analisadores;
a Seção~\ref{sec:result} discute os resultados; e a Seção~\ref{sec:concl}
conclui o trabalho.

% ======================================================================
\section{Fundamentação Teórica}\label{sec:fund}

\subsection{Lexemas, tokens e padrões}
Um \emph{lexema} é uma sequência concreta de caracteres do código-fonte, como
\texttt{nota1} ou \texttt{>=}. Um \emph{token} é o par formado por um nome
abstrato, que indica a classe do lexema, e um atributo opcional, geralmente o
próprio lexema~\cite{aho2006}. O \emph{padrão} descreve o conjunto de lexemas
que pertencem a uma classe e é normalmente expresso por uma expressão regular.
As linguagens regulares são reconhecidas por autômatos finitos, o que torna a
análise léxica executável em tempo linear no tamanho da
entrada~\cite{hopcroft2006}.

Quando mais de um padrão casa com o mesmo ponto da entrada, duas regras
resolvem a ambiguidade: a do \emph{casamento mais longo}, que prefere o maior
lexema possível (\texttt{<=} em vez de \texttt{<}), e a da \emph{prioridade},
que, entre casamentos do mesmo tamanho, escolhe o padrão listado
primeiro~\cite{cooper2011}.

\subsection{Gramáticas livres de contexto}
A estrutura sintática de uma linguagem de programação é descrita por uma
gramática livre de contexto (GLC) $G = (N, \Sigma, P, S)$, em que $N$ é o
conjunto de não-terminais, $\Sigma$ o de terminais (aqui, os tokens), $P$ o de
produções e $S$ o símbolo inicial~\cite{hopcroft2006}. A notação EBNF
estende a BNF com operadores de repetição \texttt{\{~\}} e de opcionalidade
\texttt{[~]}, o que torna a gramática mais compacta~\cite{wirth1977}.

\subsection{Análise descendente recursiva}
Um analisador descendente recursivo associa a cada não-terminal uma função
que reconhece as suas produções, consumindo tokens da esquerda para a
direita~\cite{louden2004}. Para que a escolha da produção seja
determinística com $k$ tokens de antecipação (\textit{lookahead}), a
gramática deve ser LL($k$): sem recursão à esquerda e com conjuntos
\textsc{First} disjuntos entre alternativas~\cite{grune2008}. Essa estratégia
dispensa ferramentas geradoras, produz código legível e facilita mensagens de
erro específicas, razões pelas quais é adotada por compiladores de produção e
recomendada em contextos didáticos~\cite{nystrom2021}.

% ======================================================================
\section{A Linguagem Maré}\label{sec:lingua}

A Maré é imperativa, estruturada em blocos delimitados por chaves e
estaticamente tipada, com quatro tipos primitivos: \texttt{inteiro},
\texttt{real}, \texttt{logico} e \texttt{texto}. Três decisões de projeto
guiaram a sintaxe:

\begin{itemize}
\item \textbf{Tipo à direita do nome}: \texttt{seja x: inteiro = 1;}. A
palavra inicial (\texttt{seja} ou \texttt{fixo}) identifica a declaração com
um único token de antecipação.
\item \textbf{Operadores lógicos por extenso} (\texttt{e}, \texttt{ou},
\texttt{nao}), o que aproxima as condições da leitura em português.
\item \textbf{Blocos sempre entre chaves}, inclusive em comandos de uma só
linha, o que elimina a ambiguidade do \textit{senão pendurado}.
\end{itemize}

A Listagem~\ref{lst:exemplo} mostra um trecho do programa de demonstração.
Ele contém os recursos mínimos exigidos (variáveis, bloco
\texttt{principal}, condicionais, equações matemáticas e regras lógicas) e
também funções e laços.

\begin{lstlisting}[language=Mare,caption={Trecho do programa de demonstração \texttt{completo.mare}.},label=lst:exemplo]
fixo MEDIA_MINIMA: real = 7.0;

funcao media(p1: real, p2: real,
             peso: inteiro) -> real {
    retorne (p1 + p2 * peso) / (1 + peso);
}

principal {
    seja nota1: real = 6.5;
    seja nota2: real;
    seja faltas: inteiro = 3;
    leia(nota2);
    seja m: real = media(nota1, nota2, 2);

    # regra logica composta
    se (m >= MEDIA_MINIMA e faltas <= 10
        ou nota2 == 10.0) {
        mostre("aprovada com media ", m);
    } senao se (m >= 5.0 e nao (faltas > 10)) {
        mostre("em recuperacao");
    } senao {
        mostre("reprovada");
    }
    para k de 1 ate 10 passo 2 {
        mostre(k, " ao quadrado = ", k ^ 2);
    }
    repita {
        faltas = faltas - 1;
    } ate (faltas <= 0);
}
\end{lstlisting}

A Tabela~\ref{tab:recursos} resume os comandos disponíveis.

\begin{table}[!t]
\caption{Recursos da linguagem Maré}\label{tab:recursos}
\centering\footnotesize
\begin{tabular}{@{}l>{\ttfamily}l@{}}
\toprule
\textbf{Recurso} & \normalfont\textbf{Sintaxe} \\
\midrule
Variável        & seja x: tipo [= expr]; \\
Constante       & fixo X: tipo = expr; \\
Bloco principal & principal \{ ... \} \\
Condicional     & se (c) \{\} [senao \{\}] \\ Senão-se & senao se (c) \{\} \\
Laço pré-teste  & enquanto (c) \{ ... \} \\
Laço contado    & para i de a ate b [passo p] \{\} \\
Laço pós-teste  & repita \{ ... \} ate (c); \\
Função          & funcao f(a: tipo) [-> tipo] \{\} \\
Retorno         & retorne [expr]; \\
Entrada/Saída   & leia(x); \quad mostre(e1, e2); \\
Comentários     & \# linha \quad \#* bloco *\# \\
\bottomrule
\end{tabular}
\end{table}

% ======================================================================
\section{Analisador Léxico}\label{sec:lex}

\subsection{Classes de lexemas e tokens}
A Maré possui 50 tipos de token, organizados em dez classes de lexemas
(Tabela~\ref{tab:tokens}). Os operadores lógicos, embora sejam palavras,
formam uma classe própria, porque na gramática ocupam a posição de operadores
e não de comandos.

\begin{table}[!t]
\caption{Classes de lexemas e tokens da Maré}\label{tab:tokens}
\centering\scriptsize
\setlength{\tabcolsep}{3pt}
\begin{tabular}{@{}>{\raggedright\arraybackslash}p{1.75cm}>{\raggedright\arraybackslash}p{3.3cm}>{\raggedright\arraybackslash}p{2.75cm}@{}}
\toprule
\textbf{Classe} & \textbf{Tokens} & \textbf{Lexemas} \\
\midrule
Palavra reservada & PRINCIPAL, FUNCAO, SEJA, FIXO, SE, SENAO, ENQUANTO,
  PARA, DE, ATE, PASSO, REPITA, RETORNE, MOSTRE, LEIA
  & \texttt{principal funcao seja fixo se senao enquanto para de ate passo
  repita retorne mostre leia} \\
Tipo primitivo & T\_INTEIRO, T\_REAL, T\_LOGICO, T\_TEXTO
  & \texttt{inteiro real logico texto} \\
Identificador & ID & \texttt{nota1}, \texttt{eh\_par} \\
Literal & NUM\_INT, NUM\_REAL, TEXTO, VERDADEIRO, FALSO
  & \texttt{42}, \texttt{7.0}, \texttt{"Ana"}, \texttt{verdadeiro} \\
Op.\ aritmético & MAIS, MENOS, MULT, DIV, MOD, POT & \texttt{+ - * / \% \^{}} \\
Op.\ relacional & IGUAL, DIFERENTE, MENOR, MENOR\_IGUAL, MAIOR, MAIOR\_IGUAL
  & \texttt{== != < <= > >=} \\
Op.\ lógico & E, OU, NAO & \texttt{e ou nao} \\
Atribuição & ATRIB & \texttt{=} \\
Delimitador & ABRE\_PAR, FECHA\_PAR, ABRE\_CHAVE, FECHA\_CHAVE,
  PONTO\_VIRGULA, VIRGULA, DOIS\_PONTOS, SETA
  & \texttt{( ) \{ \} ; , : ->} \\
Fim & EOF & (fim do arquivo) \\
\bottomrule
\end{tabular}
\end{table}

\subsection{Expressões regulares}
Cada classe é definida por uma expressão regular (Tabela~\ref{tab:regex}).
Espaços e comentários são reconhecidos e descartados. Algumas regras não
produzem tokens: existem para capturar entradas inválidas e emitir
mensagens de erro precisas.

\begin{table}[!t]
\caption{Expressões regulares do scanner, em ordem de prioridade}\label{tab:regex}
\centering\scriptsize
\setlength{\tabcolsep}{3pt}
\begin{tabular}{@{}r l >{\ttfamily}l@{}}
\toprule
\textbf{\#} & \textbf{Regra} & \normalfont\textbf{Expressão regular} \\
\midrule
1  & ESPACO            & [ \textbackslash t\textbackslash r]+ \\
2  & NOVA\_LINHA       & \textbackslash n \\
3  & COMENTARIO\_BLOCO & \textbackslash\#\textbackslash*[\textbackslash s\textbackslash S]*?\textbackslash*\textbackslash\# \\
4  & COMENTARIO\_ABERTO$^\dagger$ & \textbackslash\#\textbackslash* \\
5  & COMENTARIO\_LINHA & \textbackslash\#[\^{}\textbackslash n]* \\
6  & NUM\_INVALIDO$^\dagger$ & \textbackslash d+(?:\textbackslash.\textbackslash d+)?[A-Za-z\_]\textbackslash w* \\
   &                   & |\textbackslash d+\textbackslash.(?!\textbackslash d) \\
7  & NUM\_REAL         & \textbackslash d+\textbackslash.\textbackslash d+ \\
8  & NUM\_INT          & \textbackslash d+ \\
9  & TEXTO             & "(?:\textbackslash\textbackslash.|[\^{}"\textbackslash\textbackslash\textbackslash n])*" \\
10 & TEXTO\_ABERTO$^\dagger$ & "(?:\textbackslash\textbackslash.|[\^{}"\textbackslash\textbackslash\textbackslash n])* \\
11 & ID                & [A-Za-z\_][A-Za-z0-9\_]* \\
12 & SIMBOLO           & ->|==|!=|<=|>=|[<>=+\textbackslash-*/\%\^{}()\{\};,:] \\
13 & DESCONHECIDO$^\dagger$ & . \\
\bottomrule
\multicolumn{3}{@{}l}{$^\dagger$ Regra de erro: não gera token, registra um erro léxico.}
\end{tabular}
\end{table}

\subsection{Justificativa da expressão regular dos identificadores}
O identificador é definido por
\[
\mathit{ID} = (\mathit{letra} \mid \_)\,(\mathit{letra} \mid \mathit{digito} \mid \_)^{*}
\]
com $\mathit{letra} = [\texttt{A-Za-z}]$ e $\mathit{digito} = [\texttt{0-9}]$,
o que corresponde à regex \texttt{[A-Za-z\_][A-Za-z0-9\_]*}. Cada parte da
definição tem um motivo:

\begin{itemize}
\item \textbf{Primeiro caractere sem dígito.} Um lexema que começa por dígito
é sempre número. Com isso, a primeira posição decide entre as classes
\texttt{NUM} e \texttt{ID}, e \texttt{12abc} pode ser tratado como erro
sem ambiguidade (regra 6).
\item \textbf{Sublinhado permitido.} Permite nomes compostos legíveis
(\texttt{eh\_par}) e a convenção de constantes em maiúsculas
(\texttt{MEDIA\_MINIMA}).
\item \textbf{Somente ASCII.} Letras acentuadas ficam fora do alfabeto de
identificadores, e \texttt{preço} gera erro com a dica ``identificadores não
aceitam acentos''. Assim o conjunto de caracteres válidos é o mesmo em
qualquer sistema operacional ou codificação do arquivo, e as palavras
reservadas não têm grafias duplicadas (\texttt{senao} e \texttt{senão}).
\item \textbf{Fecho de Kleene no restante.} Não há limite de tamanho, e o
casamento mais longo garante que \texttt{nota12} seja um único token.
\end{itemize}

A Fig.~\ref{fig:afd} mostra o autômato finito determinístico equivalente. O
estado $q_1$ é de aceitação, e qualquer caractere fora de
$[\texttt{A-Za-z0-9\_}]$ encerra o lexema.

\begin{figure}[!t]
\centering
\begin{tikzpicture}[>=Stealth, node distance=2.6cm, auto, every state/.style={minimum size=0.8cm}, font=\footnotesize]
\node[state, initial, initial text={}] (q0) {$q_0$};
\node[state, accepting, right=of q0] (q1) {$q_1$};
\path[->] (q0) edge node {\texttt{[A-Za-z\_]}} (q1)
          (q1) edge[loop right] node {\texttt{[A-Za-z0-9\_]}} (q1);
\end{tikzpicture}
\caption{AFD que reconhece a classe \texttt{ID}.}
\label{fig:afd}
\end{figure}

\subsection{Justificativa da tabela de palavras reservadas}
As 24 palavras reservadas, os tipos e os operadores lógicos também pertencem
à linguagem de \texttt{ID}. Em vez de escrever uma expressão regular para
cada uma, o scanner reconhece o lexema como identificador e depois consulta
uma \emph{tabela de palavras reservadas}, implementada como dicionário. Se o
lexema estiver na tabela, o token é reclassificado. Essa escolha tem três
vantagens: (i) o casamento mais longo é preservado sem regras extras, e por
isso \texttt{sejam} é identificador enquanto \texttt{seja} é reservada, um
erro comum quando se escreve \texttt{seja} como regex antes de \texttt{ID};
(ii) a consulta custa $O(1)$ em média; e (iii) incluir uma palavra nova na
linguagem exige apenas uma linha na tabela. Os símbolos seguem a mesma ideia:
uma regex única os reconhece, e a tabela \texttt{SIMBOLOS} associa cada
lexema ao seu token.

\subsection{Implementação}
O scanner foi escrito em Python sem bibliotecas de terceiros. As 13 regras da
Tabela~\ref{tab:regex} são concatenadas em uma única \emph{regex mestre}, na
forma \texttt{(?P<ESPACO>...)|(?P<NOVA\_LINHA>...)|...}, e aplicada com
\texttt{re.finditer} sobre todo o código~\cite{pyre}. Como o motor de expressões regulares
do Python testa as alternativas da esquerda para a direita, a ordem das regras
implementa a regra da prioridade. As ambiguidades da linguagem são resolvidas
por essa ordem:

\begin{itemize}
\item \texttt{NUM\_INVALIDO} antes de \texttt{NUM\_INT}: \texttt{12abc} é
reportado como erro, e não aceito como \texttt{12} seguido de \texttt{abc};
\item \texttt{NUM\_REAL} antes de \texttt{NUM\_INT}: \texttt{3.14} é um só token;
\item operadores de dois caracteres antes dos de um dentro de \texttt{SIMBOLO}:
\texttt{<=} não se divide em \texttt{<} e \texttt{=};
\item comentário de bloco (\texttt{\#*}) antes do de linha (\texttt{\#}).
\end{itemize}

O nome do grupo que casou (\texttt{match.lastgroup}) indica a classe do
lexema. Cada token guarda tipo, lexema, linha e coluna. A linha é
incrementada a cada \texttt{NOVA\_LINHA} e também pelas quebras dentro de
comentários de bloco. Os erros léxicos são acumulados em uma lista, e a
varredura continua, de modo que todos são apresentados ao usuário em uma
única execução.

\subsection{Cadeia de tokens}
A Listagem~\ref{lst:mini} mostra um programa mínimo e a
Listagem~\ref{lst:cadeia} a cadeia de tokens produzida para ele, agrupada por
linha. Tokens com atributo (\texttt{ID}, números e textos) exibem o lexema.

\begin{lstlisting}[language=Mare,caption={Programa \texttt{mini.mare}.},label=lst:mini]
principal {
    seja x: inteiro = 2 + 3 * 4;
    se (x > 10 e nao falso) {
        mostre(x);
    }
}
\end{lstlisting}

\begin{lstlisting}[style=ebnf,caption={Cadeia de tokens gerada para \texttt{mini.mare}.},label=lst:cadeia]
1 <PRINCIPAL> <ABRE_CHAVE>
2 <SEJA> <ID,x> <DOIS_PONTOS> <T_INTEIRO> <ATRIB>
  <NUM_INT,2> <MAIS> <NUM_INT,3> <MULT> <NUM_INT,4>
  <PONTO_VIRGULA>
3 <SE> <ABRE_PAR> <ID,x> <MAIOR> <NUM_INT,10> <E>
  <NAO> <FALSO> <FECHA_PAR> <ABRE_CHAVE>
4 <MOSTRE> <ABRE_PAR> <ID,x> <FECHA_PAR> <PONTO_VIRGULA>
5 <FECHA_CHAVE>
6 <FECHA_CHAVE>
7 <EOF>
\end{lstlisting}

Além da cadeia, a ferramenta imprime a \emph{tabela de lexemas e tokens}, que
lista cada lexema distinto com seu token, sua classe e o número de
ocorrências. Para o programa de demonstração completo (60 linhas), foram
reconhecidos 288 tokens: 111 delimitadores, 52 identificadores, 34 literais,
33 palavras reservadas, 16 operadores aritméticos, 15 tipos, 11 atribuições,
9 operadores relacionais, 6 operadores lógicos e o EOF.

% ======================================================================
\section{Analisador Sintático}\label{sec:sint}

\subsection{Gramática}
A Listagem~\ref{lst:ebnf} apresenta a gramática da Maré em EBNF. Os terminais
aparecem entre aspas ou em maiúsculas (tokens produzidos pelo scanner), e o
símbolo inicial é \texttt{programa}. As regras se dividem em três grupos:
estrutura (programa, funções e blocos), comandos e expressões.

\begin{lstlisting}[style=ebnf,caption={Gramática livre de contexto da Maré (EBNF).},label=lst:ebnf]
programa   = {funcao | declaracao} "principal" bloco EOF
funcao     = "funcao" ID "(" [parametros] ")"
             ["->" tipo] bloco
parametros = parametro {"," parametro}
parametro  = ID ":" tipo
tipo       = "inteiro" | "real" | "logico" | "texto"
bloco      = "{" {comando} "}"

comando    = declaracao | atribuicao | chamada ";"
           | se | enquanto | para | repita
           | retorne | mostre | leia
declaracao = "seja" ID ":" tipo ["=" expr] ";"
           | "fixo" ID ":" tipo "=" expr ";"
atribuicao = ID "=" expr ";"
se         = "se" "(" expr ")" bloco
             ["senao" (se | bloco)]
enquanto   = "enquanto" "(" expr ")" bloco
para       = "para" ID "de" expr "ate" expr
             ["passo" expr] bloco
repita     = "repita" bloco "ate" "(" expr ")" ";"
retorne    = "retorne" [expr] ";"
mostre     = "mostre" "(" argumentos ")" ";"
leia       = "leia" "(" ID ")" ";"

expr       = ou
ou         = e {"ou" e}
e          = nao {"e" nao}
nao        = "nao" nao | relacional
relacional = soma [("=="|"!="|"<"|"<="|">"|">=") soma]
soma       = termo {("+"|"-") termo}
termo      = unario {("*"|"/"|"%") unario}
unario     = "-" unario | potencia
potencia   = primario ["^" unario]
primario   = NUM_INT | NUM_REAL | TEXTO | "verdadeiro"
           | "falso" | ID | chamada | "(" expr ")"
chamada    = ID "(" [argumentos] ")"
argumentos = expr {"," expr}
\end{lstlisting}

\subsection{Precedência e associatividade}
A precedência dos operadores não é tratada por tabelas externas. Ela está
embutida na hierarquia das regras de expressão: quanto mais profunda a regra,
maior a precedência do operador que ela reconhece (Tabela~\ref{tab:prec}). Na
expressão \texttt{2 + 3 * 4}, por exemplo, a regra \texttt{soma} só consome o
\texttt{+} depois que \texttt{termo} já reconheceu \texttt{3 * 4} por inteiro,
e a multiplicação fica mais funda na árvore.

\begin{table}[!t]
\caption{Precedência (crescente) e associatividade}\label{tab:prec}
\centering\footnotesize
\begin{tabular}{@{}c l >{\ttfamily}l l@{}}
\toprule
\textbf{Nível} & \textbf{Regra} & \normalfont\textbf{Operadores} & \textbf{Assoc.} \\
\midrule
1 & ou         & ou                & esquerda \\
2 & e          & e                 & esquerda \\
3 & nao        & nao               & prefixo \\
4 & relacional & == != < <= > >=   & nenhuma \\
5 & soma       & + -               & esquerda \\
6 & termo      & * / \%            & esquerda \\
7 & unario     & - (unário)        & prefixo \\
8 & potencia   & \^{}              & direita \\
\bottomrule
\end{tabular}
\end{table}

A associatividade à esquerda de \texttt{+}, \texttt{-}, \texttt{*} e
\texttt{/} vem da forma iterativa das regras (\texttt{soma = termo \{("+"|"-")
termo\}}): a cada volta do laço, a árvore acumulada torna-se o filho esquerdo
do novo nó. Assim, \texttt{10 - 3 - 2} resulta em $(10-3)-2$. A potência
é associativa à direita, como na matemática: \texttt{potencia} chama
\texttt{unario}, que pode voltar a \texttt{potencia}, e por isso
\texttt{2 \^{} 3 \^{} 2} vale $2^{9}$. Os operadores relacionais foram
definidos como não associativos. Uma cadeia como \texttt{a < b < c} gera um
erro sintático que sugere o uso de \texttt{e}, pois comparar um valor lógico
com um número não teria sentido.

\subsection{Propriedade LL e estratégia de análise}
A gramática não tem recursão à esquerda, e quase todas as alternativas de
\texttt{comando} começam por uma palavra reservada distinta. Por isso, um
token de antecipação basta para escolher a produção. A única exceção são os
comandos iniciados por \texttt{ID}, que podem ser uma atribuição
(\texttt{x = ...}) ou uma chamada (\texttt{f(...)}). Nesse ponto o parser
examina um segundo token (\texttt{=} ou \texttt{(}), o que torna essa
alternativa LL(2). O mesmo ocorre em \texttt{primario}, para distinguir uma
variável de uma chamada de função.

O parser é descendente recursivo: cada não-terminal corresponde a um método
da classe \texttt{Parser}, e a \textit{docstring} de cada método reproduz a produção
que ele implementa. A Listagem~\ref{lst:parser} mostra o método da regra
\texttt{soma}. Repetições EBNF viram laços \texttt{while}, opcionais viram
\texttt{if}, e alternativas viram testes sobre o token atual.

\begin{lstlisting}[language=Python,caption={Método da regra \texttt{soma} (simplificado).},label=lst:parser,numbers=none,xleftmargin=0.3em,framexleftmargin=0.2em]
def _soma(self):
    # soma -> termo { ("+"|"-") termo }
    no = self._termo()
    while op := self._aceitar(T.MAIS, T.MENOS):
        no = Binaria(op.lexema, no, self._termo())
    return no
\end{lstlisting}

\subsection{Integração com o analisador léxico}
Os dois analisadores se comunicam pela cadeia de tokens. O scanner produz a
lista completa, terminada em \texttt{EOF}, e o parser a percorre com um
índice. Quatro primitivas concentram o acesso: \texttt{atual} (token sob
análise), \texttt{espiar} (antecipação de $k$ tokens), \texttt{aceitar}
(consome o token se ele for do tipo esperado) e \texttt{esperar} (consome o
token obrigatório ou lança um erro sintático). Como cada token carrega linha
e coluna, os erros sintáticos são localizados com precisão no código-fonte.
Quando há erros léxicos, a análise sintática não é executada.

\subsection{Árvore sintática}
Durante o reconhecimento, cada método constrói e devolve um nó da árvore
sintática abstrata. Os nós são classes de dados (\texttt{Programa},
\texttt{Funcao}, \texttt{Declaracao}, \texttt{Se}, \texttt{Enquanto},
\texttt{Para}, \texttt{Repita}, \texttt{Binaria}, \texttt{Unaria},
\texttt{Literal}, \texttt{Variavel}, \texttt{Chamada} etc.) que expõem os
métodos \texttt{rotulo()} e \texttt{filhos()}. Isso permite que uma única
rotina desenhe a árvore no terminal e a exporte para Graphviz ou para
\LaTeX. A Fig.~\ref{fig:arvore} mostra a árvore do programa da
Listagem~\ref{lst:mini}. Nela, a precedência de \texttt{*} sobre \texttt{+}
e de \texttt{>} sobre \texttt{e} fica visível na profundidade dos nós.

\begin{figure}[!t]
\centering
\includegraphics[width=0.92\columnwidth]{arvore_mini.png}
\caption{Árvore sintática gerada pelo analisador para \texttt{mini.mare}
(exportada em formato DOT e desenhada com Graphviz).}
\label{fig:arvore}
\end{figure}

A AST difere da árvore de derivação completa porque omite os nós
intermediários sem conteúdo. Uma expressão como \texttt{x}, por exemplo, não
aparece como a cadeia \texttt{expr}$\rightarrow$\texttt{ou}$\rightarrow\dots
\rightarrow$\texttt{primario}, e sim como um único nó \texttt{Variavel}. A
AST também dispensa tokens de pontuação, cuja função se esgota na análise.
Essa forma compacta é a que será percorrida pelas fases seguintes.

\subsection{Derivação}
Para relacionar a gramática com a árvore, a derivação mais à esquerda da
expressão \texttt{2 + 3 * 4} é mostrada abaixo, em que $\Rightarrow^{*}$
resume as cadeias de regras unitárias (\texttt{relacional} $\Rightarrow$
\texttt{soma}, \texttt{unario} $\Rightarrow$ \texttt{potencia} etc.):

{\footnotesize
\begin{align*}
\textit{expr} &\Rightarrow^{*} \textit{soma}
  \Rightarrow \textit{termo}\ \texttt{+}\ \textit{termo}\\
&\Rightarrow^{*} \texttt{2}\ \texttt{+}\ \textit{termo}
  \Rightarrow \texttt{2}\ \texttt{+}\ \textit{unario}\ \texttt{*}\ \textit{unario}\\
&\Rightarrow^{*} \texttt{2 + 3}\ \texttt{*}\ \textit{unario}
  \Rightarrow^{*} \texttt{2 + 3 * 4}
\end{align*}}

\noindent Como o \texttt{*} só pode ser produzido dentro de \textit{termo},
nenhuma derivação coloca \texttt{2 + 3} como operando da multiplicação. A
gramática, portanto, não é ambígua quanto à precedência.

\subsection{Tratamento de erros}
Ao encontrar um token inesperado, o parser gera uma mensagem que informa o
que era esperado, em que contexto e o que foi encontrado, e imprime a linha
do código com um marcador sob a coluna. O ponto e vírgula esquecido recebe
um tratamento especial: se o token seguinte estiver em outra linha, o erro é
apontado no fim da linha anterior, onde o programador de fato errou.

A recuperação usa o \emph{modo pânico}~\cite{aho2006}. O erro é registrado, e
o parser descarta tokens até um ponto de sincronização, que pode ser um
\texttt{;}, um \texttt{\}} ou o início de outro comando. Um bloco
\texttt{\{...\}} encontrado nesse percurso é descartado por inteiro, com
contagem de chaves, para que o seu \texttt{\}} não feche o bloco externo por
engano. Com isso, vários erros são relatados em uma única execução, sem
cascatas de mensagens falsas.

% ======================================================================
\section{Resultados}\label{sec:result}

A implementação tem cerca de 1.150 linhas de Python, distribuídas em seis
módulos (\texttt{tokens}, \texttt{lexer}, \texttt{parser}, \texttt{arvore},
\texttt{erros} e \texttt{exibicao}), e uma interface de linha de comando que
exibe, sob demanda, as expressões regulares, a cadeia de tokens, a tabela de
lexemas e a árvore sintática. O programa de demonstração completo, de 60
linhas, que usa todos os recursos da Tabela~\ref{tab:recursos}, foi aceito, e
sua árvore reflete a precedência esperada em todas as expressões.

Para verificar a detecção de erros, foram criados dois arquivos com falhas
intencionais. No léxico, foram identificados os cinco erros inseridos: um
caractere acentuado em identificador (\texttt{preço}), um número malformado
(\texttt{12abc}), um real incompleto (\texttt{3.}), um texto sem aspas de
fechamento e um símbolo fora do alfabeto (\texttt{@}). No sintático, a
Listagem~\ref{lst:erros} reproduz os seis erros relatados em uma única
execução, cada um na posição correta.

\begin{lstlisting}[style=ebnf,caption={Saída para \texttt{erro\_sintatico.mare} (resumida).},label=lst:erros]
[2:16] esperado ':' apos o parametro 'n',
       mas foi encontrado 'inteiro'
[7:24] esperado ';' ao final da declaracao
       (faltou ';' no fim da linha)
[8:29] esperado ')' para fechar a expressao entre
       parenteses, mas foi encontrado ';'
[9:8]  esperado '(' apos 'se', mas foi encontrado 'a'
[12:5] 'senao' sem um 'se' correspondente
[15:21] comparacoes nao podem ser encadeadas;
        use 'e' (ex.: a < b e b < c)
\end{lstlisting}

\subsection{Rastreamento de um caso rejeitado}
Para mostrar o percurso completo de código, tokens e reconhecimento, a
ferramenta tem a opção \texttt{-{}-rastro}, que imprime cada regra aplicada
($\blacktriangleright$), cada token consumido (\checkmark) e cada erro ou
descarte ($\times$). O arquivo \texttt{rejeitado.mare} é o mesmo
\texttt{mini.mare} da Listagem~\ref{lst:mini}, com um único erro: falta o
\texttt{)} que fecha a condição, na linha~4
(\texttt{se (x > 10 e nao falso \{}). A cadeia de tokens dessa linha é

\smallskip
{\scriptsize\ttfamily <SE> <ABRE\_PAR> <ID,x> <MAIOR> <NUM\_INT,10> <E> <NAO> <FALSO> <ABRE\_CHAVE>}
\smallskip

\noindent e a Listagem~\ref{lst:rastro} traz o trecho final do rastreamento.

\begin{lstlisting}[style=ebnf,caption={Trecho do rastreamento de \texttt{rejeitado.mare}.},label=lst:rastro]
> se     [proximo: <SE,'se'>]
  v consome <SE,'se'>          (4:5)
  v consome <ABRE_PAR,'('>     (4:8)
  > expr  [proximo: <ID,'x'>]
    ...   (x > 10 e nao falso reconhecido)
    v consome <FALSO,'falso'>  (4:22)
  x ERRO [4:28] esperado ')' para fechar a
    condicao do 'se', mas foi encontrado '{'
  x descarta <ABRE_CHAVE,'{'> ... <FECHA_CHAVE,'}'>
v consome <FECHA_CHAVE,'}'>    (7:1)
\end{lstlisting}

A rejeição se explica pela gramática. Na regra
\texttt{se = "se" "(" expr ")" bloco}, o parser consome \texttt{se} e
\texttt{(} e chama \texttt{expr}, que reconhece
\texttt{x > 10 e nao falso} por inteiro. Ao retornar, o próximo token é
\texttt{ABRE\_CHAVE}. Nenhuma regra de expressão aceita esse token como
continuação: \texttt{\{} não é operador, e por isso \texttt{expr} termina.
A produção de \texttt{se}, por sua vez, exige \texttt{FECHA\_PAR} nessa
posição. A chamada \texttt{esperar(FECHA\_PAR)} falha, e o erro é emitido na
coluna do \texttt{\{}. A recuperação descarta o bloco \texttt{\{...\}} com as
chaves balanceadas, e o \texttt{\}} da linha~7 fecha corretamente o bloco
\texttt{principal}. Desse modo, apenas um erro é relatado, sem mensagens em
cascata.

\subsection{Testes automatizados}
Por fim, 28 testes automatizados (\texttt{unittest}) cobrem a distinção entre
palavras reservadas e identificadores, a prioridade entre as expressões
regulares, a contagem de linhas e colunas, a precedência e a associatividade
dos operadores, todos os comandos da linguagem, as mensagens de erro e o
número de erros relatados após a recuperação. Todos os testes passam.

% ======================================================================
\section{Conclusão}\label{sec:concl}

Este trabalho apresentou a especificação da linguagem Maré e a implementação
dos seus analisadores léxico e sintático. Descrever os tokens por expressões
regulares ordenadas resolveu as ambiguidades léxicas de forma simples e
verificável. Uma gramática LL projetada com cuidado, por sua vez, permitiu um
parser descendente recursivo curto e legível, em que a precedência dos
operadores é consequência direta da estrutura das regras. A opção por uma
implementação sem geradores exigiu mais código, mas deu à equipe controle
total sobre as mensagens de erro e sobre a recuperação em modo pânico.

Nas próximas etapas, a AST será a entrada da análise semântica, com tabela de
símbolos, verificação de escopo e checagem de tipos (por exemplo, rejeitar
\texttt{seja x: inteiro = "a";}). Em seguida, será construído um
interpretador ou gerador de código, o que completará o protótipo de
compilador da Maré.

O código-fonte está disponível em
\url{https://github.com/gustavoteixeirabittencourt/mare}.

% ======================================================================
\begin{thebibliography}{00}
\bibitem{aho2006} A. V. Aho, M. S. Lam, R. Sethi e J. D. Ullman,
\emph{Compilers: Principles, Techniques, and Tools}, 2.\ ed. Boston, MA,
EUA: Addison-Wesley, 2006.

\bibitem{cooper2011} K. D. Cooper e L. Torczon, \emph{Engineering a
Compiler}, 2.\ ed. Burlington, MA, EUA: Morgan Kaufmann, 2011.

\bibitem{hopcroft2006} J. E. Hopcroft, R. Motwani e J. D. Ullman,
\emph{Introduction to Automata Theory, Languages, and Computation}, 3.\ ed.
Boston, MA, EUA: Addison-Wesley, 2006.

\bibitem{louden2004} K. C. Louden, \emph{Compiladores: Princípios e
Práticas}. São Paulo, Brasil: Thomson Learning, 2004.

\bibitem{grune2008} D. Grune e C. J. H. Jacobs, \emph{Parsing Techniques: A
Practical Guide}, 2.\ ed. Nova York, NY, EUA: Springer, 2008.

\bibitem{wirth1977} N. Wirth, ``What can we do about the unnecessary
diversity of notation for syntactic definitions?'' \emph{Communications of
the ACM}, vol.~20, n.~11, pp.~822--823, 1977.

\bibitem{nystrom2021} R. Nystrom, \emph{Crafting Interpreters}. Genever
Benning, 2021.

\bibitem{pyre} Python Software Foundation, ``re --- Regular expression
operations,'' \emph{Python 3 Documentation}. [Online]. Disponível em:
\url{https://docs.python.org/3/library/re.html}
\end{thebibliography}

\end{document}
